% =============================================================
%  TubeMethodPaper_SI_body.tex
%  Copernicus (Geochronology) manuscript body, ported from thesis
%  Chapter 1 on 2026-10-30 by Paper/build_paper_from_chapter1.py.
%  Every wording change relative to the thesis is listed in
%  Paper/PORT_CHANGELOG.md. This file is now the source of truth for
%  the paper; edit it directly.
% =============================================================

\section{Ampoules and Analytical Details}\label{sec:s-analytical}

\begin{table}[t]
\centering
\caption{Silica Glass tubes that were the subject of helium step heat experiments in this work. The ampoule labeled ``Quad'' had evolved 4He quantified by isotope dilution using a Pfeiffer Prisma QMS 200 quadrupole mass spectrometer, whereas ampoules labeled ``MC'' had both evolved 3He and 4He quantified using a Helix-MC-Plus\textsuperscript{10K} multi-collector mass spectrometer. Where 3He was measured, the ampoule was soaked in tank helium with an isotope ratio of \textasciitilde0.15 Ra.}
\label{tab:ampoule-dimensions}
\resizebox{\textwidth}{!}{%
\begin{tabular}{p{1.8cm}p{1.4cm}p{1.8cm}p{1.8cm}p{1.8cm}p{1.2cm}p{1.8cm}p{2.2cm}}
\tophline
Ampoule Contents & Mass Spec. & Inner \newline Diameter \newline µm & Outer \newline Diameter \newline µm & Geometric factor \emph{G} \newline cm\textsuperscript{-2} & Volume Ratio $\left(\frac{V_{glass}}{V_{gas}}\right)$ & Content \newline Mass & Helium \newline Source \\
\middlehline
Empty & MC & 747 & 1159 & 3264 & 1.41 & --- & Tank He \\
Sulfur & MC & 529--550 & 835--872 & 5992 & ~1.5 & \textasciitilde0.2 mg & Tank He \\
Water & MC & 487--515 & 734--776 & 7789 & 1.27 & \textasciitilde0.5 mg  & Tank He \\
Durango \newline Apatite & Quad & 850 & 1350 & 2394 & 1.52 & 3.18 mg & Radiogenic \\
\bottomhline
\end{tabular}%
}
\end{table}

Heating steps and isotopic measurements for all samples were automated using custom LabView routines with the exception of the tank helium-soaked empty tube experiment on the Helix-MC-Plus\textsuperscript{10K} multi-collector mass spectrometer after the 7\textsuperscript{th} measurement. For all other ampoules, a Watlow temperature controller held the ampoule at the temperature set point for a pre-programmed duration, after which the ampoule was allowed to cool while the evolved gas was prepared for measurement for about 25 minutes. After the first 7 measurements, the empty tube experiment on the Helix-MC-Plus\textsuperscript{10K} was instead set to an elevated temperature manually while the evolved gas was allowed to accumulate continuously in an attempt to mitigate uncertainties related to frequently changing temperatures \citep{thyer2026thesis} and continued permeation at room temperature during gas cleanup and measurement. The accumulated gas was measured every 40-120 minutes. For this experiment, the temperature was changed either during the 1 minute pump-out of the lamp cell after inletting gas into the mass spectrometer, or during a period of helium accumulation for measurement. When the temperature was changed during an accumulation period, no permeability was calculated for that step, and the measured abundances of those steps were used only to track the helium contents of the tube.

Helium isotope measurements of evolved gases from step heats of one empty silica glass ampoule and glass-encapsulated sulfur and water ampoules equilibrated with 2-20 torr of high purity tank helium (R $\approx$ 0.15 R\textsubscript{a}) were acquired at Caltech on a Thermo Fisher Helix-MC-Plus\textsuperscript{10K} mass spectrometer in peak hopping mode. Helium was purified by exposing evolved gases to a charcoal trap held in liquid nitrogen, a titanium sponge getter, then a SAES NP10 getter before freely expanding into the Helix-MC-Plus\textsuperscript{10K}. 

Helium-4 evolved during the step heat of the silica glass-encapsulated Durango apatite was quantified using isotope dilution with a high purity 3He spike on a Pfeiffer Prisma QMS 200 quadrupole mass spectrometer. Evolved gas from the ampoule in the lamp heating cell was allowed to mix with the 3He spike and the mixture was frozen on charcoal held at 14~K within a cryostat, then helium was selectively released to the mass spectrometer by raising the temperature to 36~K.

Standard gasses were routinely measured to monitor drift in the sensitivities of 3He and 4He. Blanks are approximately 5 pccSTP \textsuperscript{4}He, which is very small compared to all interpreted \textsuperscript{4}He measurements. Further analytical details are given by \citet{thyer2026thesis}.

\section{Relating the Rate Constant to Glass Permeability and Ampoule Geometry}\label{sec:s-lambda-K}

Following \citet{rogers1954} and \citet{altemose1961}, steady-state permeation of helium through a glass membrane of area \emph{A} and thickness \emph{l} follows from Fick's law with a Henrian solubility \emph{S} \citep{shelby1976}; with the permeability \(K = DS\), the atom flux \emph{F} through the membrane is (full derivation in \citealp{thyer2026thesis}):

\begin{equation}
F = \  - \frac{KA\mathrm{\Delta}P}{l}
\label{eq:s1-4}
\end{equation}

Recalling that F is an atomic flux in units of \(\frac{atoms}{s}\) that
could be notated as \(\frac{dn}{dt}\):

\begin{equation}
\frac{dn}{dt} = \  - \frac{KA\mathrm{\Delta}P}{l}
\label{eq:s1-5}
\end{equation}

We assume that the helium partial pressure in the ultra-high vacuum (UHV)
heating cell surrounding the ampoule is 0, so that the pressure gradient
equals the partial pressure within the ampoule
(\(\mathrm{\Delta}P = P\)). This is a valid assumption for this work
because the heating cell is 4-5 orders of magnitude more voluminous than
the glass ampoule.

Helium is well-approximated as an ideal gas at experimental conditions, so the ideal gas law (PV = nRT) relates the pressure, \emph{P}, to the atomic abundance, \emph{n}, yielding Eq.~\eqref{eq:s1-6}:

\begin{equation}
\frac{dn}{dt} = \  - \frac{AKnRT}{lV}
\label{eq:s1-6}
\end{equation}

where \emph{R} is the gas constant \(\left( 1.36 \times 10^{- 22}\frac{cm^{3}atm}{K\ atoms} \right)\), \emph{T} is the temperature in Kelvin, and \emph{V} is the internal volume of the ampoule. Defining a geometric factor \(G = A/(lV)\) in units of cm\textsuperscript{-2} gives

\begin{equation}
\label{eq:permodeK}
\frac{dn}{dt} = \  - GKnRT
\end{equation}

Where \emph{n} is the number of helium atoms, \emph{t} is the time, \emph{G} is the geometric parameter in units of cm\textsuperscript{-2}, which is a function only of the geometry of the glass membrane, \emph{K} is the permeation constant for helium in fused silica, which follows the Arrhenius law: \(\left( K(T) = K_{0}e^{- \frac{E_{a}}{RT}} \right)\) and has been repeatedly measured for fused silica \citep{norton1953,altemose1961,swets1961,masaryk1968}, \emph{R} is the gas constant, and \emph{T} is the temperature in Kelvin. Equation~\eqref{eq:permodeK} is identical to the permeation equation used in early studies \citep{norton1953,altemose1961,swets1961,studt1962,masaryk1968}, except that it is formulated in terms of absolute atomic abundances instead of helium gas pressure by the ideal gas law. It is easily integrated by separation of variables to yield \(n_{t} = n_{0}e^{- GKRTt}\). Comparison with main-text Eq.~\eqref{eq:permode} gives \(\lambda = KGRT\).

For non-planar membranes, the effective area of permeation, \emph{A}, is
not straightforward to evaluate because the surface area on each side of
the membrane differs. \citet{masaryk1968} gives the effective area of
permeation for a cylinder (such as for this experiment) as:

\begin{equation}
A = \frac{2\pi hl}{\ln\left( \frac{b}{a} \right)}
\label{eq:s2-1}
\end{equation}

Where \emph{a} is the cylinder's internal radius, \emph{b} is the
cylinder's external radius, and \emph{h} is its length. This effective area of
permeation (Eq.~\eqref{eq:s2-1}) can be used with the equation for the internal
volume of the cylinder: \(V = h\pi a^{2}\) to evaluate the geometric
factor \emph{G} for a cylinder in terms of the ampoule
geometry. Reducing yields:

\begin{equation}
\label{eq:geomfactor}
G = \frac{2}{a^{2}\ln\left( \frac{b}{a} \right)}
\end{equation}

Interestingly, \emph{h} cancels so that the permeation rate has no
dependence on the ampoule length when permeation through the top and
bottom is negligible. Eq.~\eqref{eq:geomfactor} displays the intuitive behavior that
thinner walls and narrower tubes lead to more efficient permeation.

\section{A Numerical Model for Diffusive Permeation of Helium Through
Silica Glass
Ampoules}\label{a-numerical-model-for-diffusive-permeation-of-helium-through-silica-glass-ampoules}

Most of the permeation measurements performed on the Helix-MC-Plus\textsuperscript{10K} mass spectrometer were complicated by additional permeation occurring at variable temperature during long cooling times following high temperature experimental steps. Such a complication can be effectively avoided by limiting the duration of time for which the diffusion cell is opened to the rest of the vacuum line, or by keeping the sample at-temperature during this equilibration time. To address this, a numerical diffusive permeation model is used to calculate the ampoule's expected permeation fraction (n\textsubscript{t}/n\textsubscript{0}) for each step with a complex thermal history.

We numerically solve the radial diffusion equation for cylinders using the Forward-Time Central-Space (FTCS) Finite Difference
Method (FDM) along a radial transect of a model ampoule of arbitrary
geometry for small, discrete time intervals. The system is spatially
discretized into 100 cells along a radial transect, and numerical
stability was ensured by solving for the time interval of a single step,
$\Delta t$, such that \(\frac{D\Delta t}{\Delta x^{2}} = 0.5\), satisfying the
Courant-Friedrichs-Lewy (CFL) condition.

Helium concentrations along the transect (except for the boundaries,
where r = a and r = b) are updated according to the following equation,
an FTCS discretization of the radially symmetric diffusion equation for
a cylinder:

\[\frac{dC_{r}}{dt} = D\left\lbrack \frac{\partial^{2}C}{\partial r^{2}} + \frac{1}{r}\frac{\partial C}{\partial r} \right\rbrack\]

\begin{equation}
\Delta C_{r} \approx \Delta t\  \bullet D\left\lbrack \frac{C_{r - \Delta r} - 2C_{r} + C_{r + \Delta r}}{\Delta r^{2}} + \frac{1}{r}\frac{C_{r + \Delta r} - C_{r - \Delta r}}{2\Delta r}\  \right\rbrack
\label{eq:s5-0}
\end{equation}

Dirichlet boundary conditions were enforced at the inner and outer
boundaries. The concentration at the model ampoule's outer boundary
(where r = b) was held at 0, while the concentration at the inner
boundary (where r = a) was calculated at each discrete step by
multiplying the solubility (in units of
\(\frac{atoms}{atm \times cm^{3}}\)) by the helium partial pressure
within the ampoule, according to Eq.~\eqref{eq:s5-1}.

\begin{equation}
C_{He}^{glass} = S(T) \bullet P_{He}^{gas}
\label{eq:s5-1}
\end{equation}

In diffusive steady state, these boundary conditions give the concentration profile across the ampoule wall \citep{thyer2026thesis}:

\begin{equation}
C(r) = \ S(T) \bullet P_{He}^{gas} \bullet \frac{\ln\left( \frac{r}{b} \right)}{\ln\left( \frac{a}{b} \right)}
\label{eq:s4-3}
\end{equation}

It is sometimes useful to normalize the helium concentration of the
glass to the concentration of helium in the gas. To do this, the
internal helium partial pressure is related to an atomic or molar
concentration using the ideal gas law:

\[
C_{He}^{gas} = \frac{n}{V} = \frac{P_{He}^{gas}}{RT}
\]

The pressure term vanishes in the ratio of the two concentrations and
depends only on temperature:

\begin{equation}
C_{He}^{glass/gas} = \frac{C_{He}^{glass}}{C_{He}^{gas}} = S(T)\  \bullet RT
\label{eq:s5-2}
\end{equation}

Eq.~\eqref{eq:s5-2} strongly favors the gas and is only a weak function of temperature, varying between 2.5 and 2.8\% between 0 and 300°C.

Fluxes between the glass and the gas inside and outside the ampoule are evaluated with finite differences at the inner and outer walls and used to update the helium contents of the ampoule and of the vacuum at each time step \citep{thyer2026thesis}.

\subsubsection{Modeling permeation for ampoules that experienced complex
thermal
histories}\label{modeling-permeation-for-ampoules-that-experienced-complex-thermal-histories}

In order to compare the experimental results for ampoules that
experienced complex cooling histories with the expectations of the first
order permeation model developed in this work, the numerical model
developed above was used. The model predicts the relative helium
permeation for the three experiments on the
Helix-MC-Plus\textsuperscript{10K} multi-collector mass spectrometer for
both 4He and 3He by simulating helium permeation of a model ampoule with
the same dimensions and thermal history of the real experiment. For
ampoules with multiple measurements of the geometric dimensions, the
smaller set of measurements is used for the numerical model.

To simulate permeation during cooling of an ampoule, the temperature is
decreased at every step following a parabolic rate law determined using
the Watlow temperature records of the cooling ampoule. These functions and
the thermal data to which they are fit are reported in the dataset (\url{https://doi.org/10.5281/zenodo.23872881}). 
Further details are given by \citet{thyer2026thesis}.

\section{Derivation of the Bounding Method}\label{estimate-permeation-lag-uncertainty-for-uncorrected-long-heating-steps}

The uncertainty of the helium evolution of the encapsulated natural
material during a heating step (\(\omega\)) due to permeation lag
correction can be estimated by calculating the maximum possible helium
residing in the tube between steps, i.e. for the values for
\emph{n\textsubscript{0}} and \emph{n\textsubscript{t}}: \(n_{0}^{UB}\)
and \(n_{t}^{UB}.\ \)From the permeation lag equation (main text Eq.~\eqref{eq:permlag}), upper bound estimates for
\emph{n\textsubscript{0}} and \emph{n\textsubscript{t}} bound the
encapsulated material contribution \(\omega\) within:

\begin{equation}
m - n_{0}^{UB} \leq \omega \leq m + n_{t}^{UB}
\label{eq:s6-1}
\end{equation}

Upper bounds for \emph{n\textsubscript{t}} for a given high temperature
experimental step can be calculated if step heats are long enough to
assume steady state permeation, assuming \(\lambda\) is known or can be
conservatively estimated. To proceed, one must make the generally
conservative assumption that the production rate of helium within the
ampoule due to degassing of the encapsulated natural sample
(\(\Omega = \frac{\omega}{t}\)) is constant over time at a given
temperature. Eq.~\eqref{eq:s6-3}, given
below, shows that the amount of helium trapped in the tube (\(n_{t}\))
is proportional to the rate of helium outgassing from the natural
encapsulated sample (\(\Omega\)), so the approximation that this rate is
constant will overestimate the amount of unanalyzed helium remaining in
the ampoule after a high temperature experimental step. A new first order differential equation like Eq.~\eqref{eq:permodeK} is formulated, now
incorporating a constant source term (\(\Omega\)) to conservatively
approximate the helium influx into the ampoule from the degassing
natural sample.

\begin{equation}
\frac{dn}{dt} = \Omega - \lambda n
\label{eq:s6-2}
\end{equation}

Where \(\Omega = \frac{\omega}{t}\) and \(\lambda = KGRT\). Immediately,
the steady state quantity of helium atoms (\(n_{ss}\)) present in the
ampoule can be found by solving Eq.~\eqref{eq:s6-2} for the steady state condition
where \(\frac{dn}{dt} = 0\).

\begin{equation}
n_{ss} = \frac{\Omega}{\lambda} = \frac{\omega}{\lambda t}
\label{eq:s6-3}
\end{equation}

Additionally, the helium concentration within the ampoule during the
approach to the steady state concentration can be solved for by
integrating Eq.~\eqref{eq:s6-2} with respect to \emph{t}. First, the integrating
factor for Eq.~\eqref{eq:s6-2} is identified: \(I = e^{\lambda t}\), which is
multiplied to both sides to yield Eq.~\eqref{eq:s6-4}.

\begin{equation}
e^{\lambda t}\frac{dn}{dt} + e^{\lambda t}\lambda n = e^{\lambda t}\Omega\
\label{eq:s6-4}
\end{equation}

Then both sides are integrated with respect to \emph{t:}

\[\int_{}^{}\left( e^{\lambda t}\frac{dn}{dt} + e^{\lambda t}\lambda n \right)dt = \int_{}^{}{e^{\lambda t}\Omega}dt\]

Because
\({\frac{d}{dt}(ne^{\lambda t}) = e}^{\lambda t}\frac{dn}{dt} + e^{\lambda t}\lambda n\)
by the product rule Eq.~\eqref{eq:s6-4} reduces to:

\begin{equation}
n_{t}e^{\lambda t} = \frac{\Omega}{\lambda}e^{\lambda t} + C
\label{eq:s6-5}
\end{equation}

Where \emph{C} is an unknown constant introduced by indefinite
integration. This expression becomes more intuitive when both sides are
divided by the integrating factor and by substituting the steady state
concentration from Eq.~\eqref{eq:s6-3}.

\begin{equation}
n_{t} = n_{ss} + Ce^{- \lambda t}
\label{eq:s6-6}
\end{equation}

The integration constant \emph{C} can be solved by enforcing the
condition that \(n_{t}(t = 0) = n_{0}\), leading to:

\begin{equation}
n_{0} = n_{ss} + C
\label{eq:s6-7}
\end{equation}

Solving Eq.~\eqref{eq:s6-7} for \emph{C} and substituting into Eq.~\eqref{eq:s6-6} yields
Eq.~\eqref{eq:s6-8}. Now an inequality is assigned to the equation to reflect the
initial assumption that the isothermal helium outgassing rate
(\(\Omega\)) is constant \emph{or decreasing}.

\begin{equation}
n_{t} \leq \ n_{ss} + \left( n_{0} - n_{ss} \right)e^{- \lambda t}
\label{eq:s6-8}
\end{equation}

For long timescales (i.e. when \(e^{- \lambda t} \approx 0\)), Eq.~\eqref{eq:s6-8}
reduces to \(n_{t} \leq \ n_{ss}\). This is effectively true for steps
long enough such that \(t \geq \frac{5}{\lambda}\), for which
\(e^{- \lambda t} \leq 0.007.\) It is advisable to only consider the
bounding method for steps of sufficiently long duration such that
\(t \geq \frac{5}{\lambda}\) so that the estimation for \(n_{t}\) can be
reduced to simply the steady state value and not be dependent upon the
estimation for \(n_{0}\) or its uncertainties.

To estimate the upper bound on the permeation lag (\(n_{t}^{UB}\)), one
must consider the circumstances that would produce the highest possible
value for \(n_{t}\). In addition, the constraint that
\(t \geq \frac{5}{\lambda}\) is enforced so that previous heating steps
are not significant, simplifying the problem significantly. Because
\(n_{t}\) scales with $\omega$ by Eq.~\eqref{eq:s6-3}, the upper bound estimate
\(n_{t}^{UB}\) must assume the upper bound estimate for $\omega$, which is
given in Eq.~\eqref{eq:s6-1} as \({m + n}_{t}^{UB}\). Substituting
\(\omega^{UB} = {m + n}_{t}^{UB}\) into Eq.~\eqref{eq:s6-3} to express
\(n_{t}^{UB}\) in terms of the measured value \emph{m}:

\[
n_{t}^{UB} = n_{ss}(\omega^{UB} = m + n_{t}^{UB}) = \frac{m + n_{t}^{UB}}{\lambda t}
\]

Solving for \(n_{t}^{UB}\) yields Eq.~\eqref{eq:6}:

\begin{equation}
n_{t}^{UB} = \frac{m}{\lambda t - 1}
\label{eq:6}
\end{equation}

\section{Comparison to Literature Permeability Values}\label{comparison-to-literature-permeability-values}

Predicted and measured fractional yields for all permeation steps in ampoules besides the encapsulated water experiment agree well (within 1$\sigma$: Figs.~\ref{fig:he4-error-histogram}, \ref{fig:he3-error-histogram}) with the predictive equations presented in this work and with the permeabilities reported by \citet{swets1961} and \citet{jones1953} for all studied temperatures (24-300°C) and step fractional yields (up to \textasciitilde80\%). One outlier measurement on the empty tube experiment on the Helix-MC-Plus\textsuperscript{10K} multi-collector mass spectrometer was excluded from Figs.~\ref{fig:he4-error-histogram} and \ref{fig:he3-error-histogram}.

\begin{figure}[htbp]
\centering
\includegraphics[width=8.3cm]{Histograms4He_MC.pdf}
\caption{Histograms showing the accuracy of 4He $n_0$ predictions using permeabilities of \citet{swets1961} from either main-text Eq.~\eqref{eq:permsol} with \(\lambda = KGRT\) or the numerical diffusion model from the measured value, with shaded regions behind the histograms corresponding to 1 and 2 sigma ($\sigma$) from the average relative difference between measured and predicted values. The measured $n_0$ values of the empty and sulfur-filled ampoules are within 1$\sigma$ of the predicted values, suggesting that the permeation model and literature values apply well to those systems, and such models can be used to reliably recover helium contents of permeating ampoules for use in permeation lag corrections. On the other hand, measured and predicted 4He $n_0$ for water-filled ampoules are well resolved beyond 2$\sigma$, suggesting that their helium permeability differs strongly from the literature values. Internal reproducibilities (1$\sigma$) of a gas standard repeatedly analyzed during these experiments were better than 1\% for 4He, requiring that most of the scatter in these data be related to factors outside of analytical precision, such as long/inconsistent ramp times during heating or uncertainties or rounding of heating times.}
\label{fig:he4-error-histogram}
\end{figure}

\begin{figure}[htbp]
\centering
\includegraphics[width=8.3cm]{Histograms3He.pdf}
\caption{Similar to Fig.~\ref{fig:he4-error-histogram}, but with 3He permeabilities taken from \citet{jones1953} and \citet{swets1961}. Shaded regions behind the histograms correspond to 1 and 2 sigma ($\sigma$) from the average relative difference between the predicted and measured values. The 3He permeation model and literature permeability values apply well to empty and sulfur-filled ampoules, and such models can be used to reliably recover 3He contents of permeating ampoules for use in permeation lag corrections. On the other hand, measured and predicted values of 3He $n_0$ for water-filled ampoules are well resolved beyond 2$\sigma$, suggesting that their permeability differs strongly from the literature values. Internal reproducibilities (1$\sigma$) of a gas standard repeatedly analyzed during these experiments were better than 3\% for 3He, requiring that much of the scatter in these data be related to factors outside of analytical precision such as long/inconsistent ramp times during heating or uncertainties or rounding of heating times.}
\label{fig:he3-error-histogram}
\end{figure}

\section{Additional Analytical Corrections for the Empty Ampoule Calibration Experiments}\label{additional-analytical-corrections}

\subsection{Anomalous Signal Recalculation for the Tank Helium-Soaked Empty Ampoule}\label{anomalous-signal-recalculation}

The last high temperature step of the empty ampoule experiment on the Helix-MC-Plus\textsuperscript{10K} multi-collector mass spectrometer had an unusual and non-linear signal evolution at the beginning of the measurement for both 3He and 4He. For this measurement, 4He was recalculated using a linear regression on the last three of six measurement cycles with an R\textsuperscript{2} of 0.985 to regress to the 4He content at time 0. Typically, the signal evolution for a measurement like this would have a negative slope because the helium in the mass spectrometer is consumed as it is ionized and measured, but the slope for this measurement was positive. 3He from this measurement was calculated based on the recalculated 4He and an assumed isotope ratio for the measurement of 0.07 Ra. Three more measurements were taken over the next few days to ensure complete degassing of the ampoule. Two of these measurements yielded helium above blank, but were small compared to the final high temperature step. The yields of these steps were all considered in the calculation of the experiment's total yield.

\subsection{Independent Verification of the Encapsulated Durango Apatite Mass}\label{durango-mass-verification}

The total yield for the encapsulated Durango apatite experiment (main text, \S\ref{analytical-methods}) is estimated from the precisely-measured mass of the encapsulated apatite (3.18 mg) and a helium concentration of 194.2 nccSTP/mg \citep{farley2000}. This concentration is independently supported by the measurement of 8 different \textgreater0.5 mg aliquots of the same Durango apatite sample in-lab, which yield a concentration of 194 ± 9 nccSTP/mg.

\section{Additional Figures and Tables Referenced in the Main Text}\label{sec:s-moved}

\begin{figure}[htbp]
\centering
\includegraphics[width=8.3cm]{DurangoStepYields_RelativeAnomaly.pdf}
\caption{Measured-to-true yield ratio, \(m/\omega\), for the first 13 experimental steps (through the isothermal 300\textdegree C triplicate) of the real encapsulated Durango apatite step heat experiment (main text, \S\ref{experimental-methods}), plotted in measurement order and labeled by nominal step temperature. Here \emph{m} is the raw measured helium yield for a step and $\omega$ is its permeation lag-corrected yield; bars are anchored at a ratio of 1 (no anomaly) on a logarithmic axis, so a doubling and a halving of the ratio are given equal visual weight. Steps taken while ramping to higher temperatures are generally undermeasured (blue), as some of the sample's evolved helium has not yet fully permeated out of the tube at the time of measurement, while steps taken while ramping to lower temperatures are generally overmeasured (red), as helium retained in the tube from earlier, hotter steps continues to permeate out. The first two steps (100-150\textdegree C) are strongly overmeasured despite being on the heating ramp because little helium had yet evolved from the apatite, so measured yield there is dominated by pre-existing atmospheric helium draining from the tube rather than fresh sample degassing.}
\label{fig:durango-stepwise-yields}
\end{figure}

% AUTO-GENERATED by Chapter1/Figures/fig_insitu_lambda.py -- do not edit.
\begin{table*}[t]
\caption{In-situ permeation rate constants of the empty ampoule, fit to consecutive isothermal permeation steps with Eq.~(\ref{eq:incyield}). The first steps of the 125 and 200\,\textdegree C runs are excluded (see main text), leaving exactly determined two-step fits. The last two 250\,\textdegree C accumulations (steps 28--29) are also excluded. $n_0$ (true) is known independently from the experiment's total yield but is not used in the fits. $\lambda_{\mathrm{pred}}$ is predicted from the measured ampoule geometry and literature permeabilities \citep{swets1961,jones1953} and is shown only as a check. Uncertainties are 1$\sigma$ from the fit covariance.}
\label{tab:s-insitu-lambda}
\centering
\begin{tabular}{lrrrrrrr}
\tophline
Isotope & $T$ (\textdegree C) & Steps & $\lambda$ ($10^{-5}$\,s$^{-1}$) & $\lambda/\lambda_{\mathrm{pred}}$ & $n_0$ fit / true & $\lambda t_{\mathrm{cum}}$ & $\lambda(^3\mathrm{He})/\lambda(^4\mathrm{He})$ \\
\middlehline
$^4$He & 125 & 2 (8--9) & 2.063 $\pm$ 0.235 & 0.945 & 1.036 & 0.13 &  \\
$^3$He & 125 & 2 (8--9) & 2.075 $\pm$ 0.472 & 0.873 & 1.110 & 0.13 & 1.006 $\pm$ 0.256 \\
$^4$He & 150 & 3 (11--13) & 3.110 $\pm$ 0.074 & 0.930 & 1.033 & 0.40 &  \\
$^3$He & 150 & 3 (11--13) & 3.525 $\pm$ 0.139 & 0.964 & 1.014 & 0.45 & 1.133 $\pm$ 0.052 \\
$^4$He & 175 & 4 (14--17) & 4.662 $\pm$ 0.040 & 0.951 & 0.999 & 0.59 &  \\
$^3$He & 175 & 4 (14--17) & 5.273 $\pm$ 0.108 & 0.980 & 0.981 & 0.67 & 1.131 $\pm$ 0.025 \\
$^4$He & 200 & 2 (20--21) & 5.825 $\pm$ 0.063 & 0.841 & 1.122 & 0.44 &  \\
$^3$He & 200 & 2 (20--21) & 6.277 $\pm$ 0.392 & 0.823 & 1.167 & 0.48 & 1.078 $\pm$ 0.068 \\
$^4$He & 250 & 6 (22--27) & 12.216 $\pm$ 0.013 & 0.969 & 0.996 & 2.33 &  \\
$^3$He & 250 & 6 (22--27) & 14.233 $\pm$ 0.137 & 1.020 & 0.986 & 2.72 & 1.165 $\pm$ 0.011 \\
\bottomhline
\end{tabular}
\end{table*}


\begin{figure}[htbp]
\centering
\includegraphics[width=\textwidth]{InSituLambdaFits_Grid.pdf}
\caption{Cumulative yield against cumulative time for each isothermal run of permeation steps of the empty ampoule (rows: temperature; columns: isotope), with the rate constants fit in situ using Eq.~\eqref{eq:incyield} (red dashed curves). Hollow grey points are the first steps of the 125 and 200\textdegree C runs, which are excluded because the glass had not yet reached a diffusive steady state; the red curves for those runs are fit to the remaining two steps, and the grey dashed curves are the fits that include the first step (annotated with the change in \(\lambda\) and the resulting \(\chi^2\)/dof).}
\label{fig:s-insitu-fits}
\end{figure}
